\section{The MINERvA experiment, data and simulation samples}
\label{sec:experiment}

\subsection{Detector and beam}
MINERvA~\cite{MINERvA:2013zvz} was a fine-grained scintillator-tracker
neutrino experiment on Fermilab's NuMI beamline. Its central tracker comprises
hexagonal planes of triangular polystyrene scintillator strips in three stereo
orientations, surrounded by electromagnetic and hadronic calorimetry. The
magnetized MINOS near detector, \SI{2}{m} downstream, serves as the muon
spectrometer, measuring momentum and charge for MINOS-matched tracks.
We use the \emph{medium-energy} (ME)
neutrino-mode (FHC) exposure with $\langle E_\nu\rangle\sim\SI{6}{GeV}$.
The flux is the \emph{a posteriori} constrained NuMI ME prediction of
Ref.~\cite{MINERvA:2016iqn}, combining PPFX hadron-production weights and the
neutrino--electron-scattering constraint, identical to the published
measurement's flux~\cite{MINERvA:2021owq}.

\subsection{Data and simulation samples}
The twelve ME FHC playlists (ME1A--ME1G, ME1L--ME1P) contain
$1.057\times10^{21}$ protons on target, the same exposure as
Ref.~\cite{MINERvA:2021owq}. The selection below leaves
$4\,119\,797$ reconstructed data events.

The reference simulation uses GENIE 2.12.6~\cite{Andreopoulos:2009rq} and
MINERvA Tune v1, applied as per-event weights through MAT (MINERvA Analysis
Toolkit) reweighters. These are the flux-and-CV and GENIE dial weights,
the low-recoil 2p2h enhancement fit, and the non-resonant-pion reduction.
The full \textsc{Geant4}-based~\cite{GEANT4:2002zbu} MINERvA detector
simulation underlies the production AnaTuples. The simulated sample contains
$32\,849\,103$ true signal events, POT-normalized event-by-event to the data
exposure and processed with identical reconstruction and selection.

\subsection{Signal definition and event selection}
\label{sec:selection}
The signal definition and phase space follow Ref.~\cite{MINERvA:2021owq}
exactly. Truth-level signal events are charged-current $\nu_\mu$ interactions
in the central-tracker fiducial volume with muon phase space
\begin{equation*}
  \theta_\mu < 20^{\circ},\qquad
  1.5 < \pPar < \SI{60}{GeV/c},\qquad
  \pT < \SI{4.5}{GeV/c}.
\end{equation*}
The fiducial volume spans $5980 < z < \SI{8422}{mm}$ within an
\SI{850}{mm} apothem and contains $3.2353\times10^{30}$ nucleons, the fixed
tracker-geometry constant.

Reconstructed events require a vertex in the same fiducial volume and apothem,
a muon track with $\theta_\mu<20^{\circ}$ matched to a MINOS track with a
passing fit, no dead-time overlap, and a neutrino (rather than antineutrino)
candidate. The MINOS-match
requirement is implemented as
\texttt{isMinosMatchTrack==1 \&\& minos\_trk\_is\_ok==1}.
% (the current MINERvA-101 tutorial enforces the same \texttt{isMinosMatchTrack} match; we
% additionally require a passing MINOS fit)
The predicted genuine-background rate is \SI{0.35}{\percent} of selected
events, evaluated on playlist 1A. Genuine backgrounds comprise charged-current
events of other flavors or wrong sign, neutral-current interactions, and
out-of-fiducial-vertex events. Out-of-phase-space signal (\emph{fakes}) is
treated separately. The rate has the same sub-percent scale as the 8655
(\SI{0.2}{\percent}) predicted background events in
Ref.~\cite{MINERvA:2021owq}, which counts only wrong-flavor, wrong-sign, and
neutral-current components. Exact agreement is not expected given the
definitional and selection differences; neither prediction carries an
uncertainty. Backgrounds (reconstructed events whose true interaction fails
the signal definition) and \emph{fakes} (reconstructed in phase space, true
muon outside it) are POT-scaled and subtracted from the reconstructed-data
target before unfolding (\S\ref{sec:method}).

\subsection{Per-event analysis trees}
\label{sec:evloop}
Unbinned unfolding needs per-event records rather than histograms.
A dedicated event loop
(App.~\ref{app:codebase}) built on MAT's systematic-universe machinery writes
four flat trees per playlist: all true signal events (the efficiency
denominator), reconstructed truth-matched simulated events, simulated
background, and data. Entries carry truth and reconstructed observables
($\pT$, $\pPar$, and in extended runs $\Eavail$, $q_{3}$, $W$) and the
MnvTune central-value weight. Universe-complete
% (\texttt{\_universes\_full})
variants also carry all systematic-universe weights and laterally-shifted
observables (\S\ref{sec:systematics}). Two invariants keep the result
truth-tree authoritative:
\begin{itemize}
  \item A reconstructed signal event is filled only if its
        $(\mathrm{run},\mathrm{subrun},\mathrm{event})$ key is in the truth
        denominator. Truth-only events enter as native OmniFold ``miss'' entries.
  \item Completeness is $c=1.000000$ by construction at ME~FHC
        (\num{32849103} entries; \S\ref{sec:method-xsec}).
\end{itemize}
The loop takes one global playlist flux/calibration from the first event,
so it always runs per playlist. The twelve outputs are then merged, using
a tree-size-aware merger for
% (\texttt{TTree::SetMaxTreeSize}-aware)
the more-than-\SI{100}{GB} universe dumps. The corresponding baseline
histogram event loop supplies the flux histogram, POT-weighted across
playlists; playlist-to-playlist flux differences are a
$\sim\SI{0.2}{\percent}$ effect.



\subsection{Reco-level control distributions}
\label{sec:controlplots}
Selected data exceed the POT-scaled MnvTune-v1 prediction (signal $+$
background) by $\sim\SI{12}{\percent}$ in overall normalization.
Figure~\ref{fig:control} compares all five reconstructed observables.
The data/MC ratio has a clear $\pT$ trend, rising from $0.88$ in the lowest bin to $\sim1.2$
above \SI{1}{GeV/c}. Excess and trend together reflect the truth-level normalization deficit in
\S\ref{sec:results-model} and \S\ref{sec:3d-models}. The simulation
describes the \emph{shapes} well, which is what the unfolding relies on; the
rate offset is what the data pull corrects.

The beam and acceptance set the shapes. The $\pPar$ peak at $4$--$6$~GeV/$c$
reflects the ${\langle E_\nu\rangle}\approx\SI{6}{GeV}$ NuMI medium-energy
peak; spectrometer matching removes the lowest-$\pPar$ muons. The $\pT$
spectrum peaks near $0.7$~GeV/$c$ and falls by over an order of magnitude
across the measured range. Low-recoil (quasi-elastic and $2p2h$) events
dominate the sample. Available energy $\Eavail$ concentrates below
${\sim}\SI{0.5}{GeV}$, with a long inelastic tail in the wide upper bins;
$q_{3}$ tracks this population's momentum transfer. Reconstructed $W$ rises
through the smeared quasi-elastic and $\Delta(1232)$ region before the
deep-inelastic continuum takes over above ${\sim}\SI{2}{GeV}$.
Absolute background is at the few-percent level under the muon projections
and concentrates at high recoil (highest $\Eavail$, $q_{3}$ and $W$ bins).
Beyond the $\pT$ trend noted
above, the ratio panels also show a mild enhancement in the lowest $\Eavail$ bins,
reflecting the low-recoil excess examined at truth level in
\S\ref{sec:3d-2p2h}.



\begin{figure}[H]
  \centering
  \includegraphics[width=0.99\textwidth]{control_plots}
  \caption{Reconstructed $\pT$, $\pPar$, $\Eavail$, $q_{3}$ and $W$ for the
  analysis selection. Data points overlay POT-scaled simulated signal (blue)
  stacked on background (orange); data/MC ratios appear below. Wide hadronic
  catch bins extend to \SI{100}{GeV} and are compressed at each panel's right
  edge. The gray MC-statistical band on the ratio is mostly narrower than the points.
  The $\sim\SI{12}{\percent}$ data normalization excess over MnvTune v1
  reflects the truth-level generator deficits in \S\ref{sec:3d-models}.}
  \label{fig:control}
\end{figure}

\begin{figure}[H]
  \centering
  \includegraphics[width=0.92\textwidth]{control_corner}
  \caption{Reco-level corner plot of five observables. The lower triangle
  shows pairwise POT-scaled MC (signal$+$background) densities on a logarithmic
  color scale, with wide catch bins dropped. Each panel annotates the
  simulation's weighted Pearson correlation and the corresponding data value.
  The diagonal shows 1D marginals (data points over MC). The strongest
  correlations follow from kinematics: $q_{3}$ with $\Eavail$ ($r\approx0.81$)
  and $W$ with $\Eavail$ ($r\approx0.73$) share recoil energy; $q_{3}$ with
  $\pT$ ($r\approx0.59$) shares the muon transverse kick. Both $\Eavail$ and
  $W$ are nearly uncorrelated with the muon variables. Data reproduce every
  simulated coefficient within $0.02$.}
  \label{fig:control-corner}
\end{figure}
